%==============================================================
% BAB IV  HASIL DAN PEMBAHASAN
%==============================================================

\chapter{HASIL DAN PEMBAHASAN}

%--------------------------------------------------------------
\section{Hasil Segmentasi Pembuluh Darah}
%--------------------------------------------------------------

\subsection{Kemajuan Optimasi Pipeline}

Optimasi pipeline segmentasi pembuluh darah dilakukan melalui serangkaian percobaan ablasi
yang sistematis. Konfigurasi dasar (\textit{baseline}) menggunakan filter Gabor standar dengan
dua tahap CLAHE, menghasilkan nilai Dice 0,4469 pada set pengujian. Setiap percobaan berikutnya
menguji satu perubahan utama terhadap pipeline, sehingga kontribusi masing-masing komponen dapat
diidentifikasi secara terpisah.

Dua faktor utama memberikan peningkatan terbesar terhadap nilai Dice:

\begin{itemize}
    \item \textbf{Preprocessing top-hat}: Penambahan filter \textit{top-hat}
    morfologis sebelum filter Gabor memberikan peningkatan Dice sebesar +0,0131 (dari 0,4469
    menjadi 0,4600). Filter \textit{top-hat} mempertajam struktur gelap pembuluh darah
    terhadap latar belakang sebelum filter Gabor bekerja, sehingga respons Gabor terhadap
    pembuluh darah menjadi lebih kuat.
    \item \textbf{Fusi Gabor + Meijering}: Penggantian detektor tunggal dengan
    fusi dua detektor---Gabor untuk pembuluh besar dan Meijering untuk pembuluh tipis---memberikan
    peningkatan tambahan sebesar +0,0139 hingga mencapai nilai Dice akhir 0,4739. Ini adalah
    nilai Dice tertinggi yang dicapai dalam seluruh proses optimasi.
\end{itemize}

Konfigurasi terbaik \textit{g40\_m60\_fine} dengan formula $0{,}40 \times \text{gabor\_tophat} +
0{,}60 \times \text{meijering\_fine}$ kemudian dijadikan sebagai sumber kanal pembuluh darah
dalam skenario masukan S5.

\subsection{Temuan Penting: Dice Segmentasi $\neq$ Performa Klasifikasi}

Sebuah temuan penting yang diperoleh selama proses penelitian adalah bahwa memaksimalkan nilai
Dice segmentasi pembuluh darah tidak secara otomatis meningkatkan performa klasifikasi
\textit{stage}. Ketika konfigurasi Dice-optimal (\textit{g40\_m60\_fine}, Dice = 0,4739) digunakan
sebagai masukan CNN pada percobaan awal, performa model
justru turun dibanding menggunakan peta Gabor sederhana tanpa fusi Meijering. Penurunan
terbesar terlihat pada Stage 3, dengan F1 turun dari 0,85 menjadi 0,79.

Hal ini karena fusi Meijering yang dirancang untuk memaksimalkan tumpang tindih dengan masker
pembuluh darah sebenarnya menghaluskan peta spasial yang memuat informasi khas setiap
\textit{stage}, khususnya tekstur tepi yang membedakan Stage 2 dari Stage 3. Temuan ini
menunjukkan bahwa optimasi segmentasi dan optimasi klasifikasi adalah dua tugas yang berbeda
dan tidak bisa diasumsikan searah.

%--------------------------------------------------------------
\section{Hasil Baseline Klasik}
%--------------------------------------------------------------

\subsection{Perbandingan Antar Pengklasifikasi}

Baseline klasik menggunakan 48 fitur rekayasa manual dari peta pembuluh darah dan area tepi.
Evaluasi dilakukan menggunakan validasi silang 5-\textit{fold} yang distratifikasi pada 756
citra (kelas Stage 1, Stage 2, Stage 3 dari Zhao2024). Hasil perbandingan antar pengklasifikasi
ditunjukkan pada Tabel \ref{tab:klasik_perbandingan}.

\begin{table}[H]
\centering
\caption{Perbandingan \textit{macro F1-score} antar pengklasifikasi pada baseline klasik
(3 kelas: Stage 1, Stage 2, Stage 3)}
\label{tab:klasik_perbandingan}
\begin{tabular}{l c}
\toprule
\textbf{Pengklasifikasi} & \textbf{Macro F1} \\
\midrule
SVM (kernel RBF)           & 0,4964 \\
\textbf{Random Forest}     & \textbf{0,5147} \\
Gradient Boosting          & 0,4912 \\
\midrule
Baseline hanya fitur pembuluh darah & 0,4740 \\
\bottomrule
\end{tabular}
\end{table}

\textit{Random Forest} menghasilkan \textit{macro F1-score} terbaik sebesar 0,5147 dengan
konfigurasi 400 pohon keputusan dan bobot kelas seimbang. Penambahan 12 fitur area tepi
(\textit{ridge}) ke dalam 36 fitur pembuluh darah memberikan peningkatan nyata sebesar +0,041
(dari 0,474 menjadi 0,5147).

\subsection{Pentingnya Fitur Area Tepi}

Analisis tingkat kepentingan fitur (\textit{feature importance}) pada model \textit{Random
Forest} menunjukkan bahwa empat fitur area tepi masuk ke dalam 12 besar fitur terpenting,
seperti ditunjukkan pada Tabel \ref{tab:feature_importance}.

\begin{table}[H]
\centering
\caption{Dua belas fitur terpenting pada model \textit{Random Forest}}
\label{tab:feature_importance}
\begin{tabular}{l c}
\toprule
\textbf{Fitur} & \textbf{Tingkat Kepentingan} \\
\midrule
tortuosity\_proxy         & 0,0427 \\
comp\_size\_std           & 0,0328 \\
vessel\_anisotropy        & 0,0325 \\
int\_r\_std               & 0,0286 \\
\textbf{ridge\_main\_area}     & \textbf{0,0283} \\
\textbf{ridge\_main\_len}      & \textbf{0,0274} \\
\textbf{ridge\_main\_extent}   & \textbf{0,0267} \\
soft\_std                 & 0,0266 \\
comp\_size\_max           & 0,0260 \\
glcm\_correlation\_std    & 0,0258 \\
\textbf{ridge\_resp\_p95}      & \textbf{0,0247} \\
vessel\_density           & 0,0243 \\
\bottomrule
\end{tabular}
\end{table}

\noindent Nilai-nilai ini mengonfirmasi bahwa fitur area tepi (\textit{ridge}) memberikan
kontribusi nyata dalam membedakan \textit{stage} ROP, khususnya antara Stage 2 (yang ditandai
\textit{ridge}) dan Stage lainnya.

\subsection{Batas Atas Performa Klasik: Hambatan Stage 1}

Meskipun penambahan fitur tepi meningkatkan performa secara keseluruhan, F1-score untuk Stage
1 hampir tidak bergerak dari rentang 0,30--0,31 pada seluruh konfigurasi yang dicoba. Analisis
menunjukkan bahwa filter tepi klasik hanya mampu mendeteksi sekitar 37\% piksel \textit{demarcation
line} yang sebenarnya. Garis batas yang tipis dan berkontras rendah pada Stage 1 merupakan
struktur yang paling sulit dideteksi oleh filter berbasis aturan tanpa kemampuan belajar
adaptif.

%--------------------------------------------------------------
\section{Hasil Klasifikasi dengan Jaringan Saraf Tiruan}
%--------------------------------------------------------------

\subsection{Ringkasan Hasil Eksperimen}

Eksperimen jaringan saraf tiruan dilakukan secara bertahap, dengan setiap konfigurasi
membangun di atas hasil konfigurasi sebelumnya. Tabel \ref{tab:hasil_keseluruhan} merangkum
seluruh hasil yang diperoleh. Seluruh model jaringan saraf dievaluasi menggunakan protokol
\textit{group-aware cross-validation} (\textit{StratifiedGroupKFold} 5-\textit{fold}) yang
diterapkan secara konsisten sejak awal agar seluruh citra dari mata yang sama berada pada
\textit{fold} yang sama (lihat Subbab \ref{subsec:group_aware}).

\begin{table}[H]
\centering
\caption{Ringkasan hasil eksperimen klasifikasi}
\label{tab:hasil_keseluruhan}
\begin{tabular}{l l c l}
\toprule
\textbf{Eksperimen} & \textbf{Input} & \textbf{Macro F1} & \textbf{Catatan} \\
\midrule
CNN awal (raw RGB)          & RGB 3-kanal & 0,45  & Tanpa pra-latih \\
Baseline klasik terbaik     & 48 fitur    & 0,5147 & RF, 3 kelas \\
CNN RGB (group-aware)       & RGB 3-kanal & 0,6866 & Baseline CNN \\
\textbf{CNN softmap (group-aware)} & \textbf{Softmap} & \textbf{0,7853} & \textbf{Terbaik} \\
CNN + kalibrasi Stage1      & Softmap     & 0,8018 & Kalibrasi pasca-hoc \\
\bottomrule
\end{tabular}
\end{table}

\subsection{Protokol Evaluasi \textit{Group-Aware}}
\label{subsec:group_aware}

Pemilihan protokol evaluasi merupakan keputusan metodologis yang menentukan kredibilitas
seluruh estimasi performa. Dataset ROP yang digunakan memiliki dua karakteristik yang
menuntut pembagian data dilakukan pada tingkat kelompok mata, bukan pada tingkat citra
individual:

\begin{enumerate}
    \item \textbf{Duplikat byte-identik}: Terdapat 12 berkas citra yang persis sama dalam 6
    kelompok, termasuk tiga pasang dengan label kelas yang berbeda (konflik label). Duplikat
    semacam ini dapat muncul di kedua sisi batas pelatihan--pengujian secara bersamaan.
    \item \textbf{Citra satu mata}: Beberapa citra berasal dari mata yang sama tetapi dicatat
    sebagai berkas terpisah. Analisis \textit{normalized cross-correlation} (NCC $\geq$ 0,99)
    mengidentifikasi sekitar 63 pasang citra hampir-duplikat. Pembagian data pada tingkat citra
    akan membiarkan citra dari mata yang sama masuk ke pelatihan dan pengujian sekaligus,
    sehingga model dapat mengidentifikasi mata---bukan penyakitnya.
\end{enumerate}

Untuk mengatasi kedua karakteristik tersebut, seluruh eksperimen klasifikasi jaringan saraf
menggunakan \textit{StratifiedGroupKFold} 5-\textit{fold}, di mana seluruh citra dari satu
kelompok mata ditempatkan pada \textit{fold} yang sama. Protokol ini menghasilkan estimasi
performa yang lebih jujur dan dapat dipercaya, meskipun nilainya secara absolut lebih rendah
dibanding pembagian tingkat citra. Estimasi yang konservatif namun valid
ini lebih bermakna secara klinis dibanding angka tinggi yang tidak dapat digeneralisasikan.

\subsection{Hasil Terbaik: TinyResNetV2 dengan Masukan Softmap}

Model dengan performa terbaik adalah \textit{TinyResNetV2} yang dilatih dari awal menggunakan
masukan tiga kanal softmap, dievaluasi dengan \textit{StratifiedGroupKFold} 5-\textit{fold}.
Hasil lengkap ditunjukkan pada Tabel \ref{tab:hasil_terbaik}.

\begin{table}[H]
\centering
\caption{Hasil model terbaik: \textit{TinyResNetV2} dengan masukan softmap (group-aware CV)}
\label{tab:hasil_terbaik}
\begin{tabular}{l c}
\toprule
\textbf{Metrik} & \textbf{Nilai} \\
\midrule
OOF Macro F1-score        & \textbf{0,7853} \\
OOF Akurasi               & 0,8317 \\
OOF Macro Precision       & 0,7867 \\
OOF Macro Recall          & 0,8028 \\
\midrule
F1 per-fold (5 fold) & 0,7747 / 0,7881 / 0,7597 / 0,7834 / 0,8203 \\
\bottomrule
\end{tabular}
\end{table}

Performa per-kelas ditunjukkan pada Tabel \ref{tab:performa_per_kelas}.

\begin{table}[H]
\centering
\caption{Performa per kelas model terbaik (OOF, group-aware)}
\label{tab:performa_per_kelas}
\begin{tabular}{l c c c c}
\toprule
\textbf{Kelas} & \textbf{Presisi} & \textbf{Recall} & \textbf{F1} & \textbf{Jumlah} \\
\midrule
Normal       & 0,95 & 0,90 & 0,93 & 236 \\
Stage 1      & 0,45 & 0,77 & 0,57 & 94  \\
Stage 2      & 0,70 & 0,61 & 0,65 & 165 \\
Stage 3      & 0,83 & 0,82 & 0,83 & 261 \\
Laser Scars  & 1,00 & 0,92 & 0,96 & 343 \\
\midrule
\textbf{Macro} & \textbf{0,79} & \textbf{0,80} & \textbf{0,79} & \textbf{1.099} \\
\bottomrule
\end{tabular}
\end{table}

Visualisasi matriks konfusi model terbaik ditunjukkan pada Gambar \ref{fig:confusion_matrix}.

\begin{figure}[H]
  \centering
  \includegraphics[width=0.85\linewidth]{images/confusion_matrix_champion}
  \caption{Matriks konfusi model terbaik \textit{TinyResNetV2} dengan masukan softmap dan
  evaluasi \textit{group-aware} 5-\textit{fold} pada 1.099 citra}
  \label{fig:confusion_matrix}
\end{figure}

%--------------------------------------------------------------
\section{Pembahasan}
%--------------------------------------------------------------

\subsection{Peningkatan yang Dikontrol: Representasi Softmap vs. RGB}

Percobaan ablasi yang paling bisa dipercaya dalam penelitian ini adalah perbandingan langsung
antara masukan RGB mentah dan masukan softmap pada model, data, protokol, dan hyperparameter
yang identik---satu-satunya variabel yang berbeda adalah representasi input. Hasil perbandingan
ditunjukkan pada Gambar \ref{fig:hasil_perbandingan}.

\begin{figure}[H]
  \centering
  \includegraphics[width=0.85\linewidth]{images/results_comparison_chart}
  \caption{Perbandingan \textit{macro F1-score} seluruh konfigurasi yang dievaluasi. Bar
  yang disorot menunjukkan model terbaik}
  \label{fig:hasil_perbandingan}
\end{figure}

Masukan softmap menghasilkan \textit{macro F1-score} 0,7853 dibanding 0,6866 untuk masukan
RGB pada model yang sama, menunjukkan peningkatan terkontrol sebesar \textbf{+0,0987}. Peningkatan
ini sepenuhnya dapat dikaitkan dengan perbedaan representasi input, karena tidak ada variabel
lain yang berubah. Ini membuktikan secara langsung bahwa informasi spasial pembuluh darah
dan area tepi yang dikodekan dalam softmap lebih berguna bagi jaringan saraf dibanding nilai
piksel RGB mentah.

\subsection{Hambatan pada Stage 1 dan Stage 2}

Kelas Stage 1 dan Stage 2 tetap menjadi tantangan terbesar dalam klasifikasi. Stage 1
mencapai F1-score 0,57 dengan presisi rendah (0,45) namun recall tinggi (0,77), menunjukkan
bahwa model cenderung terlalu sering memprediksi Stage 1. Banyak sampel Stage 2 yang
diprediksi sebagai Stage 1, yang merupakan \textit{adjacent-class confusion}---kesalahan yang
dapat dipahami karena Stage 1 dan Stage 2 berbeda hanya dalam tingkat ketebalan \textit{ridge}
yang secara visual sangat mirip.

Stage 2 mencapai F1-score 0,65 dengan recall yang lebih rendah (0,61), menunjukkan sebagian
sampel Stage 2 yang bocor ke Stage 1 (nilai rendah, bergabung dengan prediksi Stage 1 yang
berlebih) dan Stage 3 (nilai tinggi). Pola ini konsisten dengan karakteristik penyakit: Stage 2
adalah kondisi \textit{transisi} antara Stage 1 dan Stage 3, sehingga batas keputusannya
secara alami lebih tidak jelas.

\subsection{Pendekatan yang Tidak Efektif}

Penelitian ini secara eksplisit mendokumentasikan dua pendekatan yang telah dicoba dan terbukti
tidak efektif:

\textbf{(1) Ordinal smoothing (v3)}: Hipotesis bahwa memperlakukan \textit{stage} sebagai
variabel ordinal (Stage 1 $<$ Stage 2 $<$ Stage 3) dengan menggunakan \textit{ordinal neighbor
smoothing} (ONSCE, $\epsilon = 0{,}15$) akan mengurangi kesalahan antar-\textit{stage} yang
berdekatan. Hasilnya menunjukkan sebaliknya: pada percobaan ablasi 2$\times$2 yang terkontrol,
penambahan \textit{ordinal smoothing} ke model v2 justru menurunkan \textit{macro F1} sebesar
0,055. \textit{Recall} Stage 1 yang ditarget untuk ditingkatkan justru turun dari 0,57 menjadi
0,51. Pendekatan ini dinyatakan gagal.

\textbf{(2) Oversampling Stage 1}: Hipotesis bahwa memperbanyak sampel Stage 1 (n=94) dalam
setiap \textit{batch} menggunakan \textit{WeightedRandomSampler} dengan faktor $\times$3 akan
memperbaiki F1-score Stage 1. Hasilnya: pada dua fold yang dijalankan, performa turun dari
0,778 menjadi 0,717 dan dari 0,773 menjadi 0,684. Masalah pada Stage 1 bukan kekurangan
eksposur (\textit{recall} sudah 0,77), melainkan terlalu banyak prediksi palsu positif
(\textit{precision} hanya 0,45). Oversampling memperburuk ketidakseimbangan ini. Pendekatan
ini dinyatakan gagal.

Pelajaran dari kedua kegagalan ini adalah bahwa masalah Stage 1--Stage 2 merupakan konfusi
yang disebabkan oleh kesamaan visual antar kelas yang berdekatan secara klinis. Solusi yang
mungkin efektif adalah pendekatan yang dapat mengurangi prediksi palsu positif Stage 1 secara
spesifik, bukan pendekatan yang memperbesar sinyal Stage 1 secara global.

\subsection{Keterbatasan Penelitian}

Beberapa keterbatasan penting perlu dicatat:

\begin{itemize}
    \item \textbf{Tidak ada identitas pasien yang terverifikasi}: Pengelompokan berdasarkan NCC
    merupakan pendekatan terbaik yang tersedia mengingat tidak adanya ID pasien dalam nama
    berkas. Meskipun ini mengurangi risiko tumpang tindih antar himpunan data, hasilnya tidak
    dapat dijamin setara dengan pembagian berbasis pasien yang terverifikasi.
    \item \textbf{Dataset tunggal}: Seluruh eksperimen dilakukan pada dataset Zhao2024 dari
    satu negara dan jenis kamera yang dominan. Kemampuan generalisasi ke dataset dari kamera
    berbeda atau populasi yang berbeda belum diuji.
    \item \textbf{Tidak ada bobot pra-latih}: Keputusan untuk tidak menggunakan bobot pra-latih
    (karena kanal masukan adalah softmap, bukan RGB) berarti penelitian ini tidak dapat
    memanfaatkan pengetahuan visual umum yang telah dipelajari dari dataset berskala besar.
    Ini adalah pertukaran (\textit{trade-off}) yang disengaja dan terdokumentasi.
    \item \textbf{Batas atas Stage 1}: F1-score Stage 1 sebesar 0,57 masih jauh di bawah
    kelas lainnya. Meningkatkan performa Stage 1 memerlukan penelitian lebih lanjut, mungkin
    dengan menggunakan model segmentasi yang dapat dipelajari (\textit{learned segmenter})
    seperti U-Net \citep{Ronneberger2015} untuk menghasilkan peta tepi yang lebih akurat.
\end{itemize}
